\section{Discussion and conclusions}
\label{sec:discussion}
The certificate interface makes the implication from a rational master bound to a convex MINLP bound explicit and checkable. Its essential connection is objective-preserving feasible-set inclusion: rigorously justified cuts retain every nonlinear feasible point, master identity binds those cuts to the discrete proof, and complete rational replay establishes the transferred bound. This allows numerical search to propose evidence without making its tolerances part of the final mathematical assertion.

A first lesson is that exactness must refer to one model throughout the calculation. Exact arithmetic in the discrete proof is insufficient if curvature analysis, derivative evaluation, or master construction interpret an expression differently. The loaded-tree contract fixes the coefficients and domains to which the certificate refers. The bound therefore concerns that exact model; relating it to an earlier modeling file is a further semantic obligation.

A second lesson is that verification failures need mathematical classification. The audits identify invalid supplied linear combinations and integer disjunctions, including steps that received external acceptance. Tiny residuals and successful process status do not justify those steps. Other rejected artifacts fail sufficient domain, curvature, or cut-enclosure tests and may still represent convex models or valid cuts. A local invalid inference refutes its supplied justification, not necessarily its final proposed bound. Retaining the evidence makes these distinctions reproducible.

A third lesson is that primal feasibility supplies information that objective comparisons cannot. Original-variable audits expose violations in two saved solver points, while independently checked original-model witnesses complete exact optimality proofs in two cases. These conclusions require bounds, integrality, rows, and objective evaluation for the same model. A feasible master point or a close library reference value cannot substitute for that work.

The uniform experiment yields accepted artifacts for 203 of 289 attempted models in separate replay, compared with 198 successful producer returns under the stated protocol. Complete proofs surviving unsuccessful producer returns further show the value of separating an artifact's validity from a process outcome. The long-rational failures show that exact values must survive text conversion and reporting as well as inference arithmetic. These results establish capability and identify failure mechanisms; the incomplete admission rules, large proof files, and shared-machine timings limit broader coverage and efficiency conclusions. Two-pass replay reduces the proof rows retained in memory without bounding the number of simultaneously live rows, and the present measurements do not establish a speed advantage over other certificate methods or solvers.

The validation has complementary roles. Executable tests and exact local audits exercise concrete checking behavior; Lean verifies the mathematical contract and the soundness of typed reference checkers for discrete proofs and rational quadratic matrices. These proofs do not verify the model loader, symbolic and interval software, or the Python master matcher and VIPR checker, nor establish their correspondence to the Lean definitions. The completed result is consequently a replayable certificate method for the stated model class with an explicit trust boundary, supported by reproducible capability and reliability evidence.
